\scanpage{252}

qui résulte aussitôt des définitions et du fait que $H$ est simpl. transitif sur $\unsure{A_0|A_{\mathrm{o}}|\vec{A}_0}$. Notons aussi

\[\begin{aligned}
U_0\times T_G\times U_1&\unsure{-|\longrightarrow}G\qquad\text{est injectif}\\
(u_0,t,u_1)&\longmapsto u_0tu_1
\end{aligned}\tag*{(104)}\]

en effet, comme $T_G$ normalise $U_0,U_1$, cela se réduit aux formules

\[T_G\cap U_0=T_G\cap U_1=\{1\}\quad(98)\quad\text{et}\]

\[\underset{B_{0G}}{\underline{U_0.T_G}}\cap U_1=\{1\}\tag*{(*)}\]

Mais on a

\[\left\{\begin{aligned}
B_{0G}\cap H&=B_{0H}\quad(=U_0.T_H)\\
B_{1G}\cap H&=B_{1H}\quad(=U_1.T_H)
\end{aligned}\right.\tag*{(105)}\]

(résulte immédiatement de $T_G\cap H=T_H$), donc on se \deleted{réduit} ramène à la formule analogue dans $H$

\[U_0.T_H\cap U_1=\{1\},\]

qui est satisfaite par hypothèse.

On a un peu \unsure{lambiné|bambiné|bousculé}, sans aller droit à la construction de (64) en termes des données $(H,U_0,\sigma,T_H)$ -- je vais y \unsure{arriver|revenir}, en introduisant

\[\left\{\begin{aligned}
B'_{0\deleted{G}}&=G_{s_0}=\{g\in G\mid gs_0=s_0\}\subset B_{0G}=G_{D_0}.\\
B'_{1\deleted{G}}&=G_{t_0}=\{g\in G\mid gt_0=t\}\subset B_{1G}=G_{D_1}.
\end{aligned}\right.\tag*{(106)}\]

On trouve en fait, plus précisément

\[\left\{\begin{aligned}
B'_{0\deleted{G}}&=T_{0\deleted{G}}.U_0\\
B'_{1\deleted{G}}&=T_{1\deleted{G}}.U_1
\end{aligned}\right.\qquad\text{produits semi-directs}\tag*{(107)}\]

% archive sheet number: 251
